%! AP Statistics — Unit 7, Lesson 7.3 — Guided Notes (Answer Key)
\documentclass[10pt]{article}
\usepackage{apstats-article}
\usepackage{apstats-key}
\newcommand{\TallMath}[1]{$\displaystyle #1\rule[-1.4em]{0pt}{3.2em}$}

\begin{document}

\pageheader{Unit 7, Lesson 7.3}{Guided Notes --- Answer Key}
\namedateperiod

\begin{objectivebox}
By the end of this lesson, I will be able to:
\begin{itemize}
    \item[$\square$] \ansline{Write a correct interpretation of a confidence interval for a population mean.}
    \item[$\square$] \ansline{Explain what ``$C$\% confident'' means in terms of repeated sampling.}
    \item[$\square$] \ansline{Identify the relationships between $n$, $C$, and the width of a CI.}
    \item[$\square$] \ansline{Use a confidence interval to justify or assess a claim about a population mean.}
\end{itemize}
\end{objectivebox}

\begin{vocabbox}
\termblanklong{Interpretation of a confidence interval}
\ansline{``We are $C$\% confident that the interval $(\ell, u)$ captures the true mean [quantity] for [population].''}

\termblanklong{``$C$\% confident'' (what it really means)}
\ansline{In repeated random sampling with the same $n$, approximately $C$\% of all confidence intervals constructed this way will capture the true $\mu$. It is a statement about the procedure, not about any one interval.}

\termblanklong{Width of a confidence interval}
\ansline{$2 \times ME = 2t^* \cdot s/\sqrt{n}$; decreases as $n$ increases; increases as $C$ increases.}

\termblanklong{Justifying a claim using a CI}
\ansline{If the claimed value $\mu_0$ lies outside the CI, there is convincing evidence against the claim. If $\mu_0$ is inside the CI, the data do not provide convincing evidence against it.}
\end{vocabbox}

\begin{hookbox}
Correct interpretation: \ans{B} \quad Why is A wrong?
\ansline{A uses ``probability,'' implying $\mu$ is random. In reality $\mu$ is fixed --- the interval either captures it or it doesn't. ``Confident'' describes our trust in the procedure.}
\end{hookbox}

\begin{notesbox}{Correct Interpretation Template (UNC-4.S)}
``We are \ans{$C$}\% confident that the interval
(\ans{$\ell$},\ \ans{$u$}) \ans{[units]}
captures the true mean \ans{[quantity]} for \ans{[population]}.''

\textbf{What does ``$C$\% confident'' mean?}\\
In \ans{repeated random sampling} with the same $n$, approximately \ans{$C$}\% of
all confidence intervals constructed this way will capture the true mean $\mu$.

\textbf{Three common errors:}
\begin{enumerate}
    \item Using ``\ans{probability}'' instead of ``confident.''
    \item Failing to reference the \ans{sample context} or \ans{population} (UNC-4.S.3).
    \item Claiming the interval contains \ans{$C$}\% of the data values.
\end{enumerate}
\end{notesbox}

\begin{notesbox}{Width of a Confidence Interval (UNC-4.U)}
Width $= 2 \times ME = 2 \cdot t^* \cdot s/\sqrt{n}$

\vspace{0.1in}
\begin{tabular}{lll}
    \textbf{Change} & \textbf{Effect on Width} & \textbf{Why} \\
    \hline
    Increase $n$ & \ans{decreases} & $ME \propto 1/\sqrt{n}$; to halve width, \ans{quadruple} $n$ \\
    Increase $C$ & \ans{increases} & Larger $t^*$ is needed \\
    Increase $s$ & \ans{increases} & More variability in the data \\
\end{tabular}

Quadrupling $n$ \ans{halves} the width.
\end{notesbox}

\begin{notesbox}{Justifying a Claim Using a CI (UNC-4.T)}
\vspace{0.1in}
\begin{tabular}{ll}
    \textbf{If $\mu_0$ is outside the CI} & \ans{Convincing} evidence that $\mu \ne \mu_0$ \\
    \textbf{If $\mu_0$ is inside the CI} & \ans{Insufficient / not convincing} evidence against the claim \\
\end{tabular}

\textbf{Example:} Sodium CI $= (2615, 3065)$ mg. Claimed value: $\mu_0 = 2300$ mg.

Inside or outside? \ans{Outside (2300 $<$ 2615).}

Conclusion: \ansline{Because 2300 mg lies below the entire interval, the data provide convincing evidence that the true mean daily sodium intake for college students exceeds 2300 mg.}
\end{notesbox}

\begin{practicebox}
\textbf{Guided Practice:} 90\% CI: $(18.3, 24.7)$ min. Claimed mean: 20 min.

\begin{enumerate}
  \item Correct interpretation:
        \ansline{We are 90\% confident that the interval (18.3, 24.7) minutes captures the true mean wait time at the clinic.}
  \item Does the CI provide convincing evidence against 20 minutes?
        \ansline{No: 20 is inside the interval (18.3, 24.7), so the data do not provide convincing evidence that the true mean differs from 20 minutes.}
  \item Effect of quadrupling $n$:
        \ansline{The margin of error is halved (since $ME \propto 1/\sqrt{n}$ and $\sqrt{4n}/\sqrt{n} = 2$). The interval width is halved as well.}
\end{enumerate}
\end{practicebox}

\end{document}
